\newcommand\pgfmathsinandcos[3]{%
  \pgfmathsetmacro#1{sin(#3)}%
  \pgfmathsetmacro#2{cos(#3)}%
}
\newcommand\LongitudePlane[3][current plane]{%
  \pgfmathsinandcos\sinEl\cosEl{#2} % elevation
  \pgfmathsinandcos\sint\cost{#3} % azimuth
  \tikzset{#1/.style={cm={\cost,\sint*\sinEl,0,\cosEl,(0,0)}}}
}
\newcommand\LatitudePlane[3][current plane]{%
  \pgfmathsinandcos\sinEl\cosEl{#2} % elevation
  \pgfmathsinandcos\sint\cost{#3} % latitude
  \pgfmathsetmacro\yshift{\cosEl*\sint}
  \tikzset{#1/.style={cm={\cost,0,0,\cost*\sinEl,(0,\yshift)}}} %
}
\newcommand\DrawLongitudeCircle[2][1]{
  \LongitudePlane{\angEl}{#2}
  \tikzset{current plane/.prefix style={scale=#1}}
   % angle of "visibility"
  \pgfmathsetmacro\angVis{atan(sin(#2)*cos(\angEl)/sin(\angEl))} %
  \draw[current plane] (\angVis:1) arc (\angVis:\angVis+180:1);
  \draw[current plane,dashed] (\angVis-180:1) arc (\angVis-180:\angVis:1);
}
\newcommand\DrawLongitudeCircleORBIT[2][1]{
  \LongitudePlane{\angEl}{#2}
  \tikzset{current plane/.prefix style={scale=#1}}
   % angle of "visibility"
  \pgfmathsetmacro\angVis{atan(sin(#2)*cos(\angEl)/sin(\angEl))} %
  \draw[current plane,color=red] (\angVis-56:1) arc (\angVis-56:\angVis+236:1);
  \draw[current plane,color=red,dashed] (\angVis-180:1) arc (\angVis-180:\angVis:1);
}
\newcommand\DrawLatitudeCircle[2][1]{
  \LatitudePlane{\angEl}{#2}
  \tikzset{current plane/.prefix style={scale=#1}}
  \pgfmathsetmacro\sinVis{sin(#2)/cos(#2)*sin(\angEl)/cos(\angEl)}
  % angle of "visibility"
  \pgfmathsetmacro\angVis{asin(min(1,max(\sinVis,-1)))}
  \draw[current plane] (\angVis:1) arc (\angVis:-\angVis-180:1);
  \draw[current plane,dashed] (180-\angVis:1) arc (180-\angVis:\angVis:1);
}

\newcommand\DrawLatitudeCircleEQUATOR[2][1]{
  \LatitudePlane{\angEl}{#2}
  \tikzset{current plane/.prefix style={scale=#1}}
  \pgfmathsetmacro\sinVis{sin(#2)/cos(#2)*sin(\angEl)/cos(\angEl)}
  % angle of "visibility"
  \pgfmathsetmacro\angVis{asin(min(1,max(\sinVis,-1)))}
  \draw[current plane, dotted] (\angVis:1) arc (\angVis:-\angVis-180:1);
  \draw[current plane, dotted] (180-\angVis:1) arc (180-\angVis:\angVis:1);
}

\newcommand\DrawLatitudeCircleORBIT[3][1]{
  \LatitudePlane{\angEl}{#2}
  \tikzset{current plane/.prefix style={scale=#1}}
  \pgfmathsetmacro\sinVis{sin(#2)/cos(#2)*sin(\angEl)/cos(\angEl)}
  % angle of "visibility"
  \pgfmathsetmacro\angVis{asin(min(1,max(\sinVis,-1)))}
  \draw[current plane,color = red] (\angVis:1) arc (\angVis:-\angVis-180:1);
  \draw[current plane,color = red] (180-\angVis:1) arc (180-\angVis:\angVis:1);
}


\scalebox{1}{\begin{tikzpicture} % "THE GLOBE" showcase
\def\R{4 } % sphere radius
\def\angEl{20} % elevation angle
\def\angAz{-20} % azimuth angle
\filldraw[ball color=white] (0,0) circle (\R);
\filldraw[fill=white] (0,0) circle (\R);

\DrawLatitudeCircle[\R]{0}; % actual equator
\DrawLongitudeCircle[\R]{-130};
\DrawLatitudeCircleEQUATOR[7]{0}; % orbit equator

\def\angEl{80} % elevation angle
\def\angAz{-90} % azimuth angle
\DrawLongitudeCircleORBIT[7]{-120};

%\DrawLatitudeCircleORBIT[7]{0}{0}; % orbit equator

\def\angEl{20} % elevation angle
\def\angAz{-20} % azimuth angle

\pgfmathsetmacro\H{\R*cos(\angEl)} % distance to north pole
\coordinate (O) at (0,0);
\coordinate (N) at (0,\H);
\coordinate (S) at (0,-\H);
\draw[left] node at (0,0){O};
\draw[left] node at (0,\H){N};
\draw[left] node at (0,-\H){S};
%\draw[thick, dashed, black](N)--(S);

\node[circle,draw,black,scale=0.3] at (O) {};
\node[circle,draw,black,scale=0.3] at (N) {};
\node[circle,draw,black,scale=0.3] at (S) {};

\coordinate (EBASE) at (-1.95 * 0.57,-2.32 * 0.57);
\coordinate (OBASE) at (-1.95 ,-2.32 );
\coordinate (EQBASE) at (-2.5 ,-1 );

\draw[dotted] (EQBASE) -- (O);
\draw[dotted] (OBASE) -- (O);
\draw[] (OBASE) -- (EBASE);

\draw[thick] ([shift=(-2.9:1cm)]OBASE) arc (0:64:1cm);


\node at (-1.6,-0.9) {$\phi$};
\node at (-1.4,-2) {$\psi$};

\coordinate (SAT) at (2,4.4);
\coordinate (POS) at (1.2,1.8);

\draw[color=blue] (SAT) -- (POS);
\node[draw, color=red, fill,circle,scale=0.4] at (SAT) {};
\node[draw, fill,circle,scale=0.4] at (POS) {};

\end{tikzpicture}}


